\chapter{2011年江西卷（理）}

\section{单选题}

\begin{problem}
若 \(z = \frac{1 + 2\i}{\i}\)，则复数 \(z =\)（\quad）

\choices
    {\(-2 - \i\)}
    {\(-2 + \i\)}
    {\(2 - \i\)}
    {\(2 + \i\)}
\end{problem}

\begin{problem}
若集合 \(A = \{x \mid -1 \leqslant 2x + 1 \leqslant 3\}\)，\(B = \left\{x \left| \frac{x - 2}{x} \leqslant 0 \right.\right\}\)，则 \(A \cap B =\)（\quad）

\choices
    {\(\{x\mid -1\leqslant x < 0\}\)}
    {\(\{x\mid 0 < x\leqslant 1\}\)}
    {\(\{x\mid 0\leqslant x\leqslant 2\}\)}
    {\(\{x\mid 0\leqslant x\leqslant 1\}\)}
\end{problem}

\begin{problem}
若 \( f(x) = \frac{1}{\sqrt{\log_{\frac{1}{2}}(2x + 1)}} \)，则 \( f(x) \) 定义域为（\quad）

\choices
    {\(\left(-\frac{1}{2},0\right)\)}
    {\(\left(-\frac{1}{2},0\right]\)}
    {\(\left(-\frac{1}{2}, +\infty\right)\)}
    {\((0, +\infty)\)}
\end{problem}

\begin{problem}
若 \( f(x) = x^2 - 2x - 4\ln x \)，则 \( f'(x) > 0 \) 的解集为（\quad）

\choices
    {\((0, +\infty)\)}
    {\((-1,0)\cup (2, + \infty)\)}
    {\((2, +\infty)\)}
    {\((-1,0)\)}
\end{problem}

\begin{problem}
已知数列 \(\{a_{n}\}\) 的前 \(n\) 项和 \(S_{n}\) 满足： \(S_{n} + S_{m} = S_{n + m}\)，且 \(a_{1} = 1\)，那么 \(a_{10} =\)（\quad）

\choices
    {1}
    {9}
    {10}
    {55}
\end{problem}

\begin{problem}
变量 \(X\) 与 \(Y\) 相对应的一组数据为 \((10, 1)\)，\((11.3, 2)\)，\((11.8, 3)\)，\((12.5, 4)\)，\((13, 5)\)；变量 \(U\) 与 \(V\) 相对应的一组数据为 \((10, 5)\)，\((11.3, 4)\)，\((11.8, 3)\)，\((12.5, 2)\)，\((13, 1)\). \(r_{1}\) 表示变量 \(Y\) 与 \(X\) 之间的线性相关系数；\(r_{2}\) 表示变量 \(V\) 与 \(U\) 之间的线性相关系数，则（\quad）

\choices
    {\(r_2 < r_1 < 0\)}
    {\(0 < r_2 < r_1\)}
    {\(r_2 < 0 < r_1\)}
    {\(r_2 = r_1\)}
\end{problem}

\begin{problem}
观察下列各式： \(5^{5} = 3125\)，\(5^{6} = 15625\)，\(5^{7} = 78125\)，\(\cdots\)，则 \(5^{2011}\) 的末四位数字为（\quad）

\choices
    {3125}
    {5625}
    {0625}
    {8125}
\end{problem}

\begin{problem}
已知 \(\alpha_{1}\)，\(\alpha_{2}\)，\(\alpha_{3}\) 是三个相互平行的平面，平面 \(\alpha_{1}\)，\(\alpha_{2}\) 之间的距离为 \(d_{1}\)，平面 \(\alpha_{2}\)，\(\alpha_{3}\) 之间的距离为 \(d_{2}\). 直线 \(l\) 与 \(\alpha_{1}\)，\(\alpha_{2}\)，\(\alpha_{3}\) 分别交于 \(P_{1}\)，\(P_{2}\)，\(P_{3}\). 那么“\(P_{1}P_{2} = P_{2}P_{3}\)”是“\(d_{1} = d_{2}\)”的（\quad）

\choices
    {充分不必要条件}
    {必要不充分条件}
    {充分必要条件}
    {既不充分也不必要条件}
\end{problem}

\begin{problem}
若曲线 \(C_{1}: x^{2} + y^{2} - 2x = 0\) 与曲线 \(C_{2}: y (y - mx - m) = 0\) 有四个不同的交点，则实数 \(m\) 的取值范围是（\quad）

\choices
    {\(\left(-\frac{\sqrt{3}}{3}, \frac{\sqrt{3}}{3}\right)\)}
    {\(\left(-\frac{\sqrt{3}}{3}, 0\right) \cup \left(0, \frac{\sqrt{3}}{3}\right)\)}
    {\(\left[-\frac{\sqrt{3}}{3}, \frac{\sqrt{3}}{3}\right]\)}
    {\(\left(-\infty, -\frac{\sqrt{3}}{3}\right) \cup \left(\frac{\sqrt{3}}{3}, +\infty\right)\)}
\end{problem}

\begin{problem}
如图，一个直径为 1 的小圆沿着直径为 2 的大圆内壁逆时针方向滚动，\(M\) 和 \(N\) 是小圆的一条固定直径的两个端点. 那么，当小圆这样滚过大圆内壁一周，点 \(M\)，\(N\) 在大圆内所绘出的图形大致是（\quad）

\FigureLayoutDeclare{img/2011/jiangxi_science/q10_fig1.png}{4.8cm}{14bp 0bp 27bp 2bp}
\bitmapfigure[max width=0.34\linewidth]{img/2011/jiangxi_science/q10_fig1.png}

\choices
    {\choicebitmap[width=0.15\paperwidth]{img/2011/jiangxi_science/q10_optA.png}}
    {\choicebitmap[width=0.15\paperwidth]{img/2011/jiangxi_science/q10_optB.png}}
    {\choicebitmap[width=0.15\paperwidth]{img/2011/jiangxi_science/q10_optC.png}}
    {\choicebitmap[width=0.15\paperwidth]{img/2011/jiangxi_science/q10_optD.png}}
\end{problem}

\section{填空题}

\begin{problem}
已知 \(|a| = |b| = 2\)，\((a + 2b) \cdot (a - b) = -2\)，则 \(a\) 与 \(b\) 的夹角为\fillinblank{}.
\end{problem}

\begin{problem}
小波通过做游戏的方式来确定周末活动，他随机地往单位圆内投掷一点，若此点到圆心的距离大于 \(\frac{1}{2}\)，则周末去看电影；若此点到圆心的距离小于 \(\frac{1}{4}\)，则去打篮球；否则，在家看书. 则小波周末不在家看书的概率为\fillinblank{}.
\end{problem}

\begin{problem}
下图是某算法程序框图，则程序运行后输出的结果是\fillinblank{}.

\texfigure[max width=0.96\linewidth]{img_repaint/2011/jiangxi_science/q13_fig1.tex}
\end{problem}

\begin{problem}
若椭圆 \(\frac{x^2}{a^2} + \frac{y^2}{b^2} = 1\) 的焦点在 \(x\) 轴上，过点 \(\left(1, \frac{1}{2}\right)\) 作圆 \(x^2 + y^2 = 1\) 的切线. 切点分别为 \(A\)，\(B\)，直线 \(AB\) 恰好经过椭圆的右焦点和上顶点，则椭圆方程是\fillinblank{}.
\end{problem}

\section{选做题}

\begin{problem}
若曲线的极坐标方程为 \(\rho = 2 \sin \theta + 4 \cos \theta\)，以极点为原点，极轴为 \(x\) 轴正半轴建立直角坐标系，则该曲线的直角坐标方程为\fillinblank{}.
\end{problem}

\begin{problem}
对于实数 \(x\)，\(y\)，若 \(|x - 1| \leqslant 1\)，\(|y - 2| \leqslant 1\)，则 \(|x - 2y + 1|\) 的最大值为\fillinblank{}.
\end{problem}

\section{解答题}

\begin{problem}
某饮料公司招聘一名员工，现对其进行一项测试，以便确定工资级别. 公司准备了两种不同的饮料共 8 杯，其颜色完全相同，并且其中 4 杯为 A 饮料，另外 4 杯为 B 饮料，公司要求此员工一一品尝后，从 8 杯饮料中选出 4 杯 A 饮料. 若 4 杯都选对，则月工资定为 3500 元；若 4 杯选对 3 杯，则月工资定为 2800 元；否则月工资定为 2100 元. 令 \(X\) 表示此人选对 A 饮料的杯数. 假设此人对 A 和 B 两种饮料没有鉴别能力.

\begin{enumerate}
\item 求 \(X\) 的分布列；
\item 求此员工月工资的期望.
\end{enumerate}
\end{problem}

\begin{problem}
在 \(\triangle ABC\) 中，角 \(A\)，\(B\)，\(C\) 的对边分别是 \(a\)，\(b\)，\(c\)，已知 \(\sin C + \cos C = 1 - \sin \frac{C}{2}\).
\begin{enumerate}
\item 求 \(\sin C\) 的值；
\item 若 \(a^2 + b^2 = 4(a + b) - 8\)，求边 \(c\) 的值.
\end{enumerate}
\end{problem}

\begin{problem}
已知两个等比数列 \(\{a_{n}\}\)，\(\{b_{n}\}\)，满足 \(a_{1} = a\)（\(a > 0\)），\(b_{1} - a_{1} = 1\)，\(b_{2} - a_{2} = 2\)，\(b_{3} - a_{3} = 3\).
\begin{enumerate}
\item 若 \(a = 1\)，求数列 \(\{a_{n}\}\) 的通项公式；
\item 若数列 \(\{a_{n}\}\) 唯一，求 \(a\) 的值.
\end{enumerate}
\end{problem}

\begin{problem}
设 \( f(x) = -\frac{1}{3} x^3 + \frac{1}{2} x^2 + 2ax \).
\begin{enumerate}
\item 若 \( f(x) \) 在 \( \left(\frac{2}{3}, +\infty\right) \) 上存在单调递增区间，求 \(a\) 的取值范围；
\item 当 \(0 < a < 2\) 时，\( f(x) \) 在 \([1, 4]\) 上的最小值为 \(-\frac{16}{3}\)，求 \( f(x) \) 在该区间上的最大值.
\end{enumerate}
\end{problem}

\begin{problem}
\(P(x_0, y_0)\)（\(x_0 \neq \pm a\)）是双曲线 \(E: \frac{x^2}{a^2} - \frac{y^2}{b^2} = 1\)（\(a > 0\)，\(b > 0\)）上一点，\(M\)，\(N\) 分别是双曲线 \(E\) 的左，右顶点. 直线 \(PM\)，\(PN\) 的斜率之积为 \(\frac{1}{5}\).
\begin{enumerate}
\item 求双曲线的离心率；
\item 过双曲线 \(E\) 的右焦点且斜率为 1 的直线交双曲线于 \(A\)，\(B\) 两点，\(O\) 为坐标原点，\(C\) 为双曲线上的一点. 满足 \(\overrightarrow{OC} = \lambda \overrightarrow{OA} + \overrightarrow{OB}\)，求 \(\lambda\) 的值.
\end{enumerate}
\end{problem}

\begin{problem}
\begin{enumerate}
    \item 如图，对于任一给定的四面体 \(A_{1}A_{2}A_{3}A_{4}\)，找出依次排列的四个相互平行的平面 \(\alpha_{1}\)，\(\alpha_{2}\)，\(\alpha_{3}\)，\(\alpha_{4}\)，使得 \(A_{i} \in \alpha_{i}\) （\(i = 1\)，\(2\)，\(3\)，\(4\)），且其中每相邻两个平面间的距离都相等；
    \item 给定依次排列的四个相互平行的平面 \(\alpha_{1}\)，\(\alpha_{2}\)，\(\alpha_{3}\)，\(\alpha_{4}\)，其中每相邻两个平面间的距离都为 1，若一个正四面体 \(A_{1}A_{2}A_{3}A_{4}\) 的四个顶点满足： \(A_{i} \in \alpha_{i}\) （\(i = 1\)，\(2\)，\(3\)，\(4\)），求该正四面体 \(A_{1}A_{2}A_{3}A_{4}\) 的体积.
\end{enumerate}

\bitmapfigure[max width=0.34\linewidth]{img/2011/jiangxi_science/q21_fig1.png}
\end{problem}
